% Данный файл распространяется под лицензией CC BY 4.0. Текст лицензии размещён на https://creativecommons.org/licenses/by/4.0/
% (c) Кафедра системного программирования, 2025

\documentclass[a4paper]{article}

\usepackage[a4paper, top=8mm, bottom=8mm, left=8mm, right=8mm]{geometry}
\usepackage{booktabs}
\usepackage{polyglossia}
\setdefaultlanguage[babelshorthands=true]{russian}
\setotherlanguage{english}

\usepackage{fontspec}
\setmainfont{FreeSerif}
\newfontfamily{\russianfonttt}[Scale=0.7]{DejaVuSansMono}

\usepackage[tiny, compact]{titlesec}
\usepackage{titling}
\setlength{\droptitle}{-1cm}
\pretitle{\begin{center}\begin{bfseries}\Large}
\posttitle{\par\end{bfseries}\end{center}}
\preauthor{\begin{center}\normalsize}
\postauthor{\par\end{center}\vspace{-1.8cm}}

\usepackage{hyperref}
\usepackage{bookmark}
\usepackage{csquotes}
\usepackage{float}

\title{Сравнение и анализ оптимизаций компиляторов}
\author{Зубенин Тимофей Сергеевич}
\date{}

\begin{document}

\maketitle

\begin{flushright}
    Группа: \emph{25.Б71-мм} \\
    Кафедра: \emph{Системного программирования} \\
    Научный руководитель: \emph{инженер-исследователь Романова Зинаида Андреевна} \\
    Номер семестра практики: \emph{2}
\end{flushright}

\section{Постановка задачи}
Целью работы является проверка возможностей современных компиляторов и их оптимизационных стратегий для обоснованного выбора компилятора под конкретную задачу. Результаты работы полезны инженерам, разрабатывающим высокопроизводительное или критически надёжное программное обеспечение.
\bigskip
\section{Описание предлагаемого решения}
\noindent
Сравнение проведено на основе статьи Филиппа Н. Хислея «Генерация высококачественного кода для программ, написанных на CИ».

В статье присутствовал код, написанный на языке C, который был использован в качестве базового для последующих проверок и экспериментов. Данный код был скомпилирован с помощью нескольких компиляторов с разными флагами, после чего было выполнено сравнение по двум основным параметрам: времени выполнения программы и объёму полученных исполняемых файлов. Сгенерированного ассемблерного код был проверен на предмет наличия в нём различных видов оптимизаций, применемых компиляторами.

\bigskip
\noindent
Ниже приведены характеристики тестовой среды:
\begin{itemize}
    \item Операционная система: Ubuntu 24.04.4 LTS
    \item Архитектура: x86-64
    \item Процессор: AMD Ryzen 5 7530U с Radeon Graphics
\end{itemize}

\medskip
\noindent
Использовались компилторы:
\begin{itemize}
    \item \texttt{clang} (версия Ubuntu 18.1.3)
    \item \texttt{gcc} (13.3.0)
    \item \texttt{compcert} (3.17)
\end{itemize}

\medskip
\noindent
В ходе замеров времени и объёма исполянемого файла возникло две проблемы:

\medskip
\noindent
\textbf{Проблема 1: деление на ноль.}\\
В оригинальном коде присутствовала проверка того, как компиляторы оптимизируют деление на ноль. При компиляции выдаётся предупреждение, но сам код с делением не удаляется, что позволяет компилятору пометить часть кода как недостижимую.

\noindent
Этот блок мешал замерам времени и проверке оптимизаций, поэтому он был убран с помощью флага компиляции.

\bigskip
\noindent
\textbf{Проблема 2: выход за границы массива \texttt{invector5}.}\\
В статье проверяется, в первую очередь, удаление переменной индукции; выход за пределы массива является второстепенным, но приводит к неопределённому поведению (undefined behavior). Для чистоты экспериментов и корректных замеров массив был расширен, чтобы исключить выход за границы без влияния на проверяемую оптимизацию.

\bigskip
\bigskip
\bigskip
\bigskip

\section{Эксперименты}
\noindent
Для каждого компилятора были выполнены сборки с тремя режимами оптимизации:
\begin{itemize}
    \item \texttt{-O0} (без оптимизации);
    \item \texttt{-Os} (оптимизация по размеру кода);
    \item \texttt{-O3} (агрессивная оптимизация по скорости).
\end{itemize}
\smallskip
Измерялось среднее время работы 100 последовательных запусков (фиксированный поток) и размер исполняемого файла. Также по методике из статьи проверялось наличие различных типов оптимизаций.

\bigskip
\subsection{Результаты замеров времени и размера}
\renewcommand{\arraystretch}{1.2}
\begin{table}[H]
\centering
\caption{Результаты для \texttt{-O0}}
\begin{tabular}{lccc}
\toprule
& \textbf{ccomp} & \textbf{gcc} & \textbf{clang} \\
\midrule
Время работы (сек) & 0,00113 & 0,00114 & 0,00115 \\
Размер файла (бит) & 17072   & 16992   & 16992 \\
\bottomrule
\end{tabular}
\end{table}

\\
\begin{table}[H]
\centering
\caption{Результаты для \texttt{-Os}}
\begin{tabular}{lccc}
\toprule
& \textbf{ccomp} & \textbf{gcc} & \textbf{clang} \\
\midrule
Время работы (сек) & 0,00112 & 0,00115 & 0,00113 \\
Размер файла (бит) & 17032   & 16992   & 17048 \\
\bottomrule
\end{tabular}
\end{table}

\\
\begin{table}[H]
\centering
\caption{Результаты для \texttt{-O3}}
\begin{tabular}{lccc}
\toprule
& \textbf{ccomp} & \textbf{gcc} & \textbf{clang} \\
\midrule
Время работы (сек) & 0,00112 & 0,00112 & 0,00110 \\
Размер файла (бит) & 17032   & 16992   & 17048 \\
\bottomrule
\end{tabular}
\end{table}

\bigskip
\bigskip


Размер результирующего файла при использовании флагов \texttt{-Os} (оптимизация на размер) и \texttt{-O3} (агрессивная оптимизация по скорости) оказывается одинаковым для каждого из рассмотренных компиляторов, полученные исполняемые файлы не совпадают по своей внутренней структуре.
\medskip


Интересно, что для компилятора \texttt{clang} размер исполняемого файла при оптимизации \texttt{-Os} неожиданно оказался больше, чем при отключении оптимизации \texttt{-O0} (17048 бит против 16992). На первый взгляд это противоречит самой идее оптимизации по размеру, ведь обычно ожидается, что \texttt{-Os} уменьшает объём кода. Это объясняется тем, что при оптимизации по объёму компилятор сначала применяет некоторые преобразования, направленные на увеличение скорости выполнения, и только затем пытается сократить размер. В данном конкретном случае эти преобразования, повышающие скорость, привели к нежелательному росту объёма кода, что и дало такой необычный результат для \texttt{clang}.

\medskip





\subsection{Поддерживаемые оптимизации}
В таблице приведено, какие оптимизации из рассмотренных в статье поддерживаются каждым компилятором (знак «+» означает поддержку, «-» — отсутствие, «+-» — частичную поддержку).

\begin{table}[H]
\centering
\caption{Поддержка оптимизаций компиляторами}
\begin{tabular}{lccc}
\toprule
\textbf{Оптимизация} & \textbf{clang} & \textbf{gcc} & \textbf{ccomp} \\
\midrule
Размножение копий               & + & - & + \\
Свертка констант                & + & + & + \\
Алгебраические упрощения        & + & + & + \\
Удаление излишних загрузок/сохранений & + & + & + \\
Удаление излишних присваиваний  & + & + & + \\
Снижение мощности               & + & + & – \\
Удаление циклов                 & + & + & – \\
Удаление переменных индукции циклов & + & + & – \\
Удаление общих подвыражений     & + & + & – \\
Вынесение инвариантного кода    & – & – & – \\
Удаление недостижимого кода     & + & + & +- \\
Разворачивание циклов           & – & – & – \\
Слияние циклов                  & – & – & – \\
Удаление лишних цепочек перехода & + & + & + \\
\bottomrule
\end{tabular}
\end{table}

Как видно из таблицы, \texttt{gcc} и \texttt{clang} поддерживают подавляющее большинство рассмотренных оптимизаций. \texttt{CompCert} заметно отстаёт — его применение оправдано в системах с высокими требованиями к математической корректности компиляции (например, в авионике или медицинском оборудовании), где производительность вторична.

\bigskip
\section{Заключение}
При измерении времени выполнения тестовой программы все три компилятора демонстрируют практически одинаковые результаты, которые очень улучшаются при оптимизации. Это происходит из-за малой вычислительной сложности (отсутствие тяжёлых циклов, ресурсоёмких операций ввода‑вывода или сложных математических вычислений) изначального кода, различия в производительности лежат в пределах погрешности измерений. Иными словами, для столь простой задачи любой из уровней оптимизации даёт сопоставимый результат.

\medskip
\noindent
В ходе выполнения работы получены следующие результаты:
\begin{itemize}
    \item Выполнены замеры времени выполнения и размера исполняемых файлов для трёх уровней оптимизации.
    \item Установлено, что \texttt{gcc} и \texttt{clang} обеспечивают широкий набор оптимизаций, включая удаление циклов, снижение мощности и удаление общих подвыражений, в то время как \texttt{compcert} ориентирован на формальную верификацию, а не на производительность.
    \item Практически идентичное время работы объясняется простотой тестовых примеров; для реальных приложений различия между компиляторами могут быть существенными.
\end{itemize}

\smallskip
\noindent
Хочется отметить, что
\begin{itemize}
    \item Ассемблерный код \texttt{clang} оказался наиболее легко читаемым: компилятор достаточно строго сохраняет порядок операций, что упрощает ручной анализ.
    \item Код \texttt{gcc} воспринимается заметно сложнее из-за большего разнообразия и необычности используемых инструкций, а также активного перемешивания арифметических операций с целыми и нецелыми переменными.
    \item \texttt{CompCert} генерирует наиболее запутанный код: он дополнительно помещает адрес переменной в регистр для каждого изменения значения, что сильно затрудняет понимание при наличии нескольких переменных.
\end{itemize}
\bigskip
\noindent
Таким образом, для рассмотренной платформы в проектах, где предполагается регулярный просмотр ассемблерного кода, предпочтительнее использовать \texttt{clang}. \texttt{Clang} генерирует более читаемый и структурированный ассемблерный код, что существенно упрощает его анализ и отладку.
\bigskip

\noindent
В свою очередь, \texttt{gcc} может дать большее преимущество, когда важен минимальный объём исполянемого файла. Из замеров сделанных ранее, можно заметить, что по этому показателю gcc опережает остальные рассмотренный компиляторы на файле для сравнения.
\bigskip

\noindent
Что касается \texttt{CompCert}, его следует выбирать осознанно и только в ситуациях, где требуется формально верифицированный компилятор. Несмотря на гарантии корректности трансляции, по набору оптимизаций он существенно уступает остальным, поэтому его применение оправданно лишь в задачах, где надёжность кода критичнее производительности. CompCert также выгодно отличился при работе с флагом -O0. Например, в случае с лишними цепочками переходов он справился как с -O0, так и с -O3. GCC и Clang же выполнили это только с -O3, поэтому в случае, когда нужна компиляция без агрессивной оптимизации (поскольку она затратна по ресурсам) и при условии большого количества потенциальных цепочек переходов, CompCert может с этим помочь.
\bigskip

\noindent
Материалы по данной работе модно найти по адресу: \url{https://github.com/tzubenin/practice_2sem}.

\end{document}
